% ============================================================
% Part 1 — Advanced Worked Problems
% Topics: double-angle formulae, t-formulae, products-to-sums,
%         restricting domains, multi-step proof problems
% ============================================================

% ------------------------------------------------------------------
\begin{problem}[Proving Pythagoras via the $t$-Formulae]
In a right-angled triangle, let one acute angle be $\theta$. If the hypotenuse
is $c$, the legs may be written as
\[
a=c\sin\theta,
\qquad
b=c\cos\theta.
\]
Let $t=\tan(\theta/2)$.

\begin{enumerate}[(i)]
  \item Express $\sin\theta$ and $\cos\theta$ in terms of $t$.
  \item Substitute these expressions into $a^2+b^2$ and simplify.
  \item Hence verify that $a^2+b^2=c^2$.
\end{enumerate}
\end{problem}

\begin{hintbox}
Use
$\sin\theta=\dfrac{2t}{1+t^2}$ and
$\cos\theta=\dfrac{1-t^2}{1+t^2}$. After squaring, the numerator should become
a perfect square.
\end{hintbox}

\begin{solution}
The $t$-formulae give
\[
\sin\theta=\frac{2t}{1+t^2},
\qquad
\cos\theta=\frac{1-t^2}{1+t^2}.
\]
Therefore
\[
a^2+b^2
=c^2\left[
\left(\frac{2t}{1+t^2}\right)^2+
\left(\frac{1-t^2}{1+t^2}\right)^2
\right].
\]
Combining the fractions,
\[
a^2+b^2
=c^2\left[
\frac{4t^2+(1-t^2)^2}{(1+t^2)^2}
\right]
=c^2\left[
\frac{4t^2+1-2t^2+t^4}{(1+t^2)^2}
\right].
\]
The numerator is
\[
1+2t^2+t^4=(1+t^2)^2.
\]
Hence
\[
a^2+b^2=c^2.
\]
\end{solution}

\begin{takeaways}
\begin{itemize}[leftmargin=*]
  \item The $t$-formulae rationalise sine and cosine using one algebraic
        parameter.
  \item The identity $\sin^2\theta+\cos^2\theta=1$ appears here as the algebraic
        simplification
        $4t^2+(1-t^2)^2=(1+t^2)^2$.
  \item Good substitutions should reduce trigonometry to ordinary algebra.
\end{itemize}
\end{takeaways}
% ------------------------------------------------------------------

\begin{problem}[The Cosine Sum-to-Product Identity in a Triangle]
Let $A$, $B$, and $C$ be the interior angles of a triangle, so
$A+B+C=180^\circ$.

\begin{enumerate}[(i)]
  \item Using
  \[
  \cos P+\cos Q
  =2\cos\frac{P+Q}{2}\cos\frac{P-Q}{2},
  \]
  show that
  \[
  \cos2A+\cos2B=-2\cos C\cos(A-B).
  \]
  \item Use $\cos2C=2\cos^2C-1$ to show that
  \[
  \cos2A+\cos2B+\cos2C
  =
  -2\cos C[\cos(A-B)-\cos C]-1.
  \]
  \item Hence prove that
  \[
  \cos2A+\cos2B+\cos2C
  =
  -1-4\cos A\cos B\cos C.
  \]
\end{enumerate}
\end{problem}

\begin{hintbox}
Use $A+B=180^\circ-C$, so $\cos(A+B)=-\cos C$. In the final part, simplify
$\cos(A-B)+\cos(A+B)$.
\end{hintbox}

\begin{solution}
Using the sum-to-product identity with $P=2A$ and $Q=2B$,
\[
\cos2A+\cos2B=2\cos(A+B)\cos(A-B).
\]
Since $A+B=180^\circ-C$, $\cos(A+B)=-\cos C$, and hence
\[
\cos2A+\cos2B=-2\cos C\cos(A-B).
\]

Adding $\cos2C=2\cos^2C-1$ gives
\[
\cos2A+\cos2B+\cos2C
=-2\cos C\cos(A-B)+2\cos^2C-1,
\]
so
\[
\cos2A+\cos2B+\cos2C
=-2\cos C[\cos(A-B)-\cos C]-1.
\]

Now $\cos C=\cos(180^\circ-(A+B))=-\cos(A+B)$, so the bracket is
\[
\cos(A-B)+\cos(A+B)=2\cos A\cos B.
\]
Therefore
\[
\cos2A+\cos2B+\cos2C
=-2\cos C(2\cos A\cos B)-1
=-1-4\cos A\cos B\cos C.
\]
\end{solution}

\begin{takeaways}
\begin{itemize}[leftmargin=*]
  \item Triangle identities often use $A+B=180^\circ-C$ as the key move.
  \item Choosing the right double-angle form can create a useful factor.
  \item Scaffolded proofs usually reveal the structure of the final identity.
\end{itemize}
\end{takeaways}
% ------------------------------------------------------------------

\begin{problem}[The Triple Tangent Identity and Angle Deduction]
\begin{enumerate}[(i)]
  \item Using
  \[
  \tan(X+Y)=\frac{\tan X+\tan Y}{1-\tan X\tan Y},
  \]
  prove
  \[
  \tan(A+B+C)
  =
  \frac{\tan A+\tan B+\tan C-\tan A\tan B\tan C}
       {1-\tan A\tan B-\tan B\tan C-\tan C\tan A}.
  \]
  \item In any triangle $ABC$, show that $\tan(A+B+C)=0$.
  \item Hence prove that, for any non-right-angled triangle,
  \[
  \tan A+\tan B+\tan C=\tan A\tan B\tan C.
  \]
  \item In an acute triangle, $\tan A=2$ and $\tan B=3$. Find the exact size of
        $\angle ACB$.
\end{enumerate}
\end{problem}

\begin{hintbox}
Group the angles as $(A+B)+C$. In part (iii), a fraction is zero when its
numerator is zero, provided the expression is defined.
\end{hintbox}

\begin{solution}
Let $x=\tan A$, $y=\tan B$, and $z=\tan C$. Then
\[
\tan(A+B)=\frac{x+y}{1-xy}.
\]
Applying the addition formula again,
\[
\tan(A+B+C)
=
\frac{\frac{x+y}{1-xy}+z}
     {1-\frac{x+y}{1-xy}z}.
\]
Multiplying numerator and denominator by $1-xy$ gives
\[
\tan(A+B+C)=
\frac{x+y+z-xyz}{1-xy-yz-zx},
\]
which is the required identity.

In a triangle, $A+B+C=180^\circ$, so
\[
\tan(A+B+C)=\tan180^\circ=0.
\]
For a non-right-angled triangle the tangent expressions are defined, so the
numerator must be zero:
\[
\tan A+\tan B+\tan C-\tan A\tan B\tan C=0.
\]
Thus
\[
\tan A+\tan B+\tan C=\tan A\tan B\tan C.
\]

Finally,
\[
2+3+\tan C=2\cdot3\cdot\tan C,
\]
so $5+\tan C=6\tan C$ and $\tan C=1$. Since the triangle is acute,
\[
\angle ACB=45^\circ.
\]
\end{solution}

\begin{takeaways}
\begin{itemize}[leftmargin=*]
  \item The tangent addition formula can be iterated to handle three angles.
  \item In a triangle, $A+B+C=180^\circ$ turns the triple-angle tangent into
        zero.
  \item The identity
        $\tan A+\tan B+\tan C=\tan A\tan B\tan C$ is valid only when none of the
        triangle's angles is $90^\circ$.
\end{itemize}
\end{takeaways}
% ------------------------------------------------------------------

\begin{problem}[Harmonic Addition and Equations]
\begin{enumerate}[(i)]
  \item Express $\cos x-\sin x$ in the form
        \[
        R\cos(x+\alpha),
        \]
        where $R>0$ and $0<\alpha<\pi/2$.
  \item Hence solve
        \[
        \cos x-\sin x=\frac{\sqrt2}{2},
        \qquad 0\le x\le2\pi.
        \]
\end{enumerate}
\end{problem}

\begin{hintbox}
Expand $R\cos(x+\alpha)$ and match the coefficients of $\cos x$ and $\sin x$.
\end{hintbox}

\begin{solution}
Using
\[
R\cos(x+\alpha)=R\cos x\cos\alpha-R\sin x\sin\alpha,
\]
we need
\[
R\cos\alpha=1,
\qquad
R\sin\alpha=1.
\]
Thus
\[
R=\sqrt{1^2+1^2}=\sqrt2,
\qquad
\tan\alpha=1,
\]
so $\alpha=\pi/4$. Therefore
\[
\cos x-\sin x=\sqrt2\cos\left(x+\frac{\pi}{4}\right).
\]

The equation becomes
\[
\sqrt2\cos\left(x+\frac{\pi}{4}\right)=\frac{\sqrt2}{2},
\]
so
\[
\cos\left(x+\frac{\pi}{4}\right)=\frac12.
\]
For $0\le x\le2\pi$, the argument $x+\pi/4$ lies in
$[\pi/4,9\pi/4]$. Hence
\[
x+\frac{\pi}{4}=\frac{\pi}{3},\frac{5\pi}{3},\frac{7\pi}{3}.
\]
Therefore
\[
x=\frac{\pi}{12},\frac{17\pi}{12},\frac{25\pi}{12}.
\]
Only the first two lie in $[0,2\pi]$, so
\[
\boxed{x=\frac{\pi}{12},\frac{17\pi}{12}}.
\]
\end{solution}

\begin{takeaways}
\begin{itemize}[leftmargin=*]
  \item A linear combination of $\sin x$ and $\cos x$ can be rewritten as one
        shifted sinusoid.
  \item Coefficient matching determines both the amplitude $R$ and the phase
        angle.
  \item Solve the transformed equation in the shifted argument, then convert
        back to $x$.
\end{itemize}
\end{takeaways}
% ------------------------------------------------------------------

\begin{problem}[The Symmetry of Auxiliary Angles]
For $a,b>0$, write
\[
a\cos x+b\sin x
\]
in the four forms
\[
R\sin(x+\alpha),\qquad
R\sin(x-\beta),\qquad
R\cos(x+\gamma),\qquad
R\cos(x-\delta),
\]
where $R=\sqrt{a^2+b^2}$ and
$0\le\alpha,\beta,\gamma,\delta<2\pi$.

\begin{enumerate}[(i)]
  \item For $a=b=1$, find $\alpha,\beta,\gamma,\delta$.
  \item For general $a,b>0$, show that $\alpha+\delta$ and $\beta+\gamma$ are
        constant, and find their values.
\end{enumerate}
\end{problem}

\begin{hintbox}
Expand each auxiliary-angle form and compare the coefficients of $\cos x$ and
$\sin x$. Be careful with the signs in the minus forms.
\end{hintbox}

\begin{solution}
For $a=b=1$,
\[
\cos x+\sin x=\sqrt2\sin\left(x+\frac{\pi}{4}\right),
\]
so $\alpha=\pi/4$. Also
\[
\sqrt2\sin(x-\beta)
=\sqrt2(\sin x\cos\beta-\cos x\sin\beta).
\]
Matching coefficients gives $\cos\beta=1/\sqrt2$ and
$-\sin\beta=1/\sqrt2$, hence $\beta=7\pi/4$.

Similarly,
\[
\sqrt2\cos(x+\gamma)
=\sqrt2(\cos x\cos\gamma-\sin x\sin\gamma),
\]
so $\gamma=7\pi/4$, and
\[
\sqrt2\cos(x-\delta)
=\sqrt2(\cos x\cos\delta+\sin x\sin\delta),
\]
so $\delta=\pi/4$.

For general $a,b>0$, matching coefficients gives
\[
\sin\alpha=\frac{a}{R},
\qquad
\cos\alpha=\frac{b}{R},
\]
and
\[
\cos\delta=\frac{a}{R},
\qquad
\sin\delta=\frac{b}{R}.
\]
Thus $\alpha$ and $\delta$ are complementary acute angles, so
\[
\alpha+\delta=\frac{\pi}{2}.
\]

For the two negative-shift forms, the coefficient matching places
\[
\beta=2\pi-\alpha,
\qquad
\gamma=2\pi-\delta.
\]
Therefore
\[
\beta+\gamma
=4\pi-(\alpha+\delta)
=4\pi-\frac{\pi}{2}
=\frac{7\pi}{2}.
\]
\end{solution}

\begin{takeaways}
\begin{itemize}[leftmargin=*]
  \item The four auxiliary-angle forms describe the same sinusoid from
        different phase-shift conventions.
  \item The sine and cosine forms differ by a phase shift of $\pi/2$.
  \item For $a,b>0$, complementary-angle relationships explain the constant
        sums.
\end{itemize}
\end{takeaways}
